\documentclass[11pt]{article}

\usepackage[a4paper,margin=1in]{geometry}
\usepackage{amsmath,amssymb,amsfonts}
\usepackage{booktabs}
\usepackage{url}
\usepackage[hidelinks]{hyperref}
\usepackage{enumitem}
\usepackage[numbers,sort&compress]{natbib}
\bibliographystyle{unsrtnat}

\newcommand{\doi}[1]{\href{https://doi.org/#1}{\texttt{doi:#1}}}

\title{Canonical Optimisation Problems in Biology and Bioinformatics:\\
Mathematical Formulations and Benchmarking Perspectives}

\author{}
\date{}

\begin{document}

\maketitle

\section*{Introduction}

Optimisation problems lie at the core of modern computational biology and bioinformatics.
Many foundational tasks can be expressed as well-defined mathematical optimisation problems,
often involving combinatorial search, nonconvex energy landscapes, or constrained programming.

These problems remain computational bottlenecks despite decades of algorithmic progress and
are therefore ideal candidates for systematic benchmarking of classical, AI-based, quantum,
and emerging computational paradigms. The central question is how well current heuristics
perform in practice, and where novel approaches---including deep learning, reinforcement
learning, and quantum-classical hybrids---can meaningfully improve on the state of the art.

This document presents six canonical optimisation problems:

\begin{itemize}[noitemsep]
    \item Multiple Sequence Alignment (MSA)
    \item Phylogenetic Tree Search
    \item Haplotype Phasing
    \item RNA Secondary Structure Prediction
    \item Radiotherapy Treatment Planning
    \item Metabolic Flux Optimisation (FBA)
\end{itemize}

For each problem we provide: a contextual overview of its biological and clinical relevance;
the specific computational barriers that make it a persistent bottleneck;
a precise mathematical formulation; an account of the current heuristic landscape 
and the opportunities for improvement by classical, AI-based, and quantum approaches;
a description of existing benchmark datasets; and selected references.

\subsection*{Overview of the Six Problems}

\textbf{Multiple Sequence Alignment (MSA)} is one of the most pervasive operations in
computational biology. It is usual the first step in comparative
genomics and proteomics pipeline: identifying conserved functional domains, constructing
phylogenies, predicting protein structure and function, and detecting evolutionary
signatures of selection. With the explosion of next-generation sequencing (NGS) data,
the scale of MSA problems has grown dramatically---modern pangenome analyses involve
thousands to millions of sequences. While heuristic tools such as ClustalW, MUSCLE,
and MAFFT are widely used, their greedy progressive strategies degrade in accuracy on
divergent sequences and at pangenomic scale. Improving alignment quality in these
regimes is an active open challenge.
\medskip

\textbf{Phylogenetic Tree Search} is central to evolutionary biology, epidemiology,
and public health. Accurate phylogenies underpin tracking of pathogen evolution
(as demonstrated during COVID-19 and HIV surveillance), ancestral sequence reconstruction,
and comparative genomics. The tree topology search space grows very rapidly with the number of taxa,
making exhaustive search infeasible beyond a few dozen.
State-of-the-art tools such as RAxML and IQ-TREE use hill-climbing heuristics on the
likelihood surface and are highly effective for moderate datasets, but struggle as
datasets reach thousands of taxa or when pandemic surveillance demands near-real-time
phylodynamic inference \cite{Felsenstein1981,Yang2014}.
\medskip

\textbf{Haplotype Phasing} is a critical step in personal genomics and precision medicine.
The human genome is diploid: each individual carries two copies of each chromosome.
Haplotype-resolved genomes reveal compound heterozygosity, long-range regulatory effects,
and population structure that are invisible to standard genotype calls. Despite massive
investment in sequencing, most clinical and GWAS datasets are delivered as unphased
genotypes. Existing heuristics (HapCUT2, WhatsHap) are well-suited to short reads but
struggle with coverage gaps, indels, and high heterozygosity in long-read data.
As long-read technologies (PacBio, Oxford Nanopore) become standard in clinical
sequencing, more accurate and scalable phasing methods are needed.
\medskip

\textbf{RNA Secondary Structure Prediction} is foundational to structural biology and
RNA therapeutics. RNA molecules fold into complex secondary structures that determine
their function: ribosome assembly, splice-site recognition, CRISPR guide RNA targeting,
and mRNA vaccine stability all depend critically on secondary structure.
The classical minimum free energy (MFE) model \cite{ZukerStiegler1981,Mathews1999}
has been the standard approach for over four decades, but the imprecision of empirical
energy parameters caps accuracy on standard benchmarks. Predicting RNA structure at
genome scale---particularly for long non-coding RNAs, pseudoknotted structures, and
structured viral genomes---remains a genuinely open challenge.
\medskip

\textbf{Radiotherapy Treatment Planning} is a high-stakes optimisation problem
directly affecting patient outcomes. Intensity-modulated radiotherapy (IMRT) and
volumetric modulated arc therapy (VMAT) require computing beam fluence maps that
deliver a prescribed dose to a tumour while minimising exposure to surrounding healthy
tissue and organs at risk. The problem couples a large-scale quadratic or
mixed-integer programme with inherently uncertain biological models for tumour response
and normal tissue complication. Beam angle selection, online adaptive replanning,
and biological objective functions are active bottlenecks where current heuristics
leave significant room for improvement \cite{Bortfeld2006,Craft2007}.
\medskip

\textbf{Metabolic Flux Optimisation (Flux Balance Analysis, FBA)} is the workhorse
model of systems biology. It predicts the steady-state flux distribution in a metabolic
network by solving a linear programme over stoichiometric constraints, and is used
in drug target identification, metabolic engineering, synthetic biology, and
personalised medicine. Genome-scale metabolic models (GEMs) for human cells now
include thousands of reactions and metabolites. The core LP is fast to solve, but
key challenges lie in model accuracy---particularly the specification of uptake bounds
and multi-omics integration---and in hard extensions such as mixed-integer FBA,
dynamic FBA, and ensemble modelling, all of which are active open problems
\cite{VarmaPalsson1994,Orth2010}.

\newpage

%%%%%%%%%%%%%%%%%%%%%%%%%%%%%%%%%%%%%%%%%%
\section{Multiple Sequence Alignment (MSA)}
%%%%%%%%%%%%%%%%%%%%%%%%%%%%%%%%%%%%%%%%%%

\subsection*{Problem Description}

Given $k$ biological sequences, the goal is to align them by inserting gaps so that
homologous positions are aligned and an overall scoring function is maximized.

\subsection*{Mathematical Formulation}

Let $S_1, \dots, S_k$ be sequences and $A$ an alignment configuration.

The classical sum-of-pairs formulation is:

\[
\max_{A} \sum_{1 \le i < j \le k} \text{Score}(A_i, A_j)
\]

where
\[
\text{Score}(A_i, A_j) = \sum_{p=1}^{L} s(a_{ip}, a_{jp}) + \text{gap penalties}
\]

with $L$ the alignment length.

\subsection*{Existing Benchmark Datasets}

The MSA field has a rich set of benchmark databases, all based on reference alignments
derived from experimentally solved three-dimensional protein structures.

\textbf{BAliBASE} \cite{Thompson1999,BAliBASE3} is the most widely used MSA benchmark.
Maintained at the IGBMC in Strasbourg (J.~Thompson), it provides manually curated
reference alignments organised into reference sets targeting specific alignment
difficulties: high sequence divergence, unequal N/C-terminal extensions, internal
insertions, transmembrane sequences, and structural repeats. BAliBASE~3.0 contains
hundreds of families; the evaluation metrics are the Sum-of-Pairs score (SP) and
Total Column score (TC) relative to the structural reference. 
Available at \url{https://www.igbmc.fr/bioinformatics/BAliBASE}.

\textbf{HomFam} \cite{HomFam2010} is a large-scale homology-based benchmark derived from
Pfam and containing families of diverse sizes (up to tens of thousands of sequences),
designed to stress-test aligners at pangenomic scale. It complements BAliBASE by
covering the large-$k$ regime where progressive heuristics degrade most.

\textbf{OXBench} \cite{OXBench2003} and \textbf{PREFAB} \cite{PREFAB2004} provide
complementary reference sets based on structural alignments, with different taxonomic
coverage and difficulty distributions.

\textbf{SABmark} \cite{SABmark2005} specifically targets twilight-zone alignments (very
low sequence identity, $<$25\%), representing the hardest regime for classical substitution
matrices.

\textbf{Key limitation for BENQODHI purposes:} all existing MSA benchmarks evaluate
\emph{biological correctness} (agreement with structural reference) rather than
\emph{computational hardness} in the sense required for quantum-vs-classical comparison.
There is no benchmark that exposes SP-score optima and measures approximation ratios;
this gap is precisely what a quantum-facing benchmark would need to fill.

\subsection*{Heuristic Landscape and Improvement Opportunities}

The established heuristic benchmarks are progressive aligners: ClustalW, MUSCLE, 
and MAFFT build alignments by first constructing a guide tree and then merging 
pairwise alignments in a greedy bottom-up pass. This approach is fast but makes 
irreversible decisions, causing accuracy to degrade on divergent or large datasets.

There is meaningful room for improvement, particularly in two regimes: (i) 
\emph{high divergence}, where classical substitution matrices fail to capture deep 
evolutionary signals; and (ii) \emph{pangenomic scale}, where the number of input 
sequences reaches millions and current tools run out of memory or accuracy.

AI approaches are making inroads. Transformer-based architectures operating directly 
on unaligned sequences---such as BetaAlign, which reframes alignment as a 
sequence-to-sequence translation problem---match or exceed classical tools in accuracy 
on standard benchmarks while capturing long-range coevolutionary patterns that 
window-based heuristics miss \cite{Dotan2024}. The learnMSA framework integrates 
protein language model embeddings (e.g.\ ProtT5) into a differentiable hidden Markov 
model for progressive alignment, showing accuracy gains on very divergent families 
\cite{learnMSA2024}. End-to-end learning with a differentiable Smith--Waterman layer 
has also demonstrated that jointly optimising alignment and downstream tasks improves 
both \cite{Petti2023}. The main open challenge for deep learning methods is 
\emph{generalisation}: current models often show strong intra-family performance but 
weaker inter-family performance, and scaling to thousands of sequences simultaneously 
remains an architecture-level challenge.
Quantum approaches for MSA remain at an exploratory stage, with QUBO formulations of 
pairwise sub-problems having been studied but not yet competitive at practical scales.

\subsection*{References}

\begin{itemize}[noitemsep]
  \item \cite{CarrilloLipman1988}: Branch-and-bound framework for optimal sum-of-pairs 
    MSA using pairwise projection bounds (SIAM J.\ Appl.\ Math., 1988).
    \doi{10.1137/0148063}
  \item \cite{WangJiang1994}: NP-completeness of MSA under the SP-score
    (J.\ Comput.\ Biol., 1994). \doi{10.1089/cmb.1994.1.337}
  \item \cite{Notredame2007}: Review of progressive and consistency-based MSA methods
    (PLoS Comput.\ Biol., 2007). \doi{10.1371/journal.pcbi.0030123}
  \item \cite{Thompson1999}: BAliBASE, benchmark alignment database for MSA evaluation
    (Nucleic Acids Res., 1999). \doi{10.1093/nar/27.1.599}
  \item \cite{BAliBASE3}: BAliBASE~3.0, latest benchmark release with extended reference
    sets (Proteins, 2005). \doi{10.1002/prot.20527}
  \item \cite{Dotan2024}: BetaAlign, transformer-based MSA via sequence-to-sequence 
    translation (Bioinformatics, 2025). \doi{10.1093/bioinformatics/btaf009}
  \item \cite{learnMSA2024}: learnMSA2, deep protein alignment integrating protein 
    language models (Bioinformatics, 2024). \doi{10.1093/bioinformatics/btae534}
  \item \cite{Petti2023}: End-to-end learning of MSA with differentiable Smith--Waterman
    (Bioinformatics, 2023). \doi{10.1093/bioinformatics/btac724}
  \item Wikipedia: \url{https://en.wikipedia.org/wiki/Multiple_sequence_alignment}
\end{itemize}

%%%%%%%%%%%%%%%%%%%%%%%%%%%%%%%%%%%%%%%%%%
\section{Phylogenetic Tree Search}
%%%%%%%%%%%%%%%%%%%%%%%%%%%%%%%%%%%%%%%%%%

\subsection*{Problem Description}

Given aligned sequences, infer the evolutionary tree topology and branch lengths that
best explain the data under a probabilistic model.

\subsection*{Mathematical Formulation}

Let $T$ denote tree topology and $\theta$ branch lengths/model parameters.

Maximum Likelihood formulation:

\[
\max_{T,\theta} \mathcal{L}(T,\theta \mid D)
= \prod_{i=1}^{L} P(D_i \mid T, \theta)
\]

Search space size for $n$ taxa: $(2n-5)!!$

\subsection*{Existing Benchmark Datasets}

Unlike MSA, phylogenetics does not have a single dominant curated benchmark.
Practice instead relies on a combination of empirical datasets and simulation.

\textbf{TreeBASE} \cite{TreeBASE} is a repository of published phylogenetic matrices and
trees, containing thousands of empirical datasets from the primary literature. It provides
a heterogeneous collection of real-world instances of varying size and complexity, though
ground-truth topologies are unknown for empirical data.

\textbf{Simulated benchmarks} generated with tools such as
\texttt{seq-gen} \cite{Rambaut1997} and \texttt{INDELible} \cite{Fletcher2009} allow
known-topology evaluation. The standard protocol is to simulate sequences under a given
evolutionary model on a known tree, then assess how well inference methods recover it
(measured by Robinson--Foulds distance \cite{Robinson1981} and branch score).

\textbf{CRW (Comparative RNA Web)} \cite{CRW2002} and the \textbf{1000 Genomes
phylogenetic datasets} provide large-scale empirical alignments used in recent
benchmarking of scalable tools such as RAxML-NG and FastTree.

\textbf{The COVID-19 phylodynamic datasets} (GISAID \cite{GISAID}) have served as
de facto benchmarks for near-real-time inference at scale ($>$10$^5$ sequences),
exposing the limits of current hill-climbing approaches.

\textbf{Key limitation for BENQODHI purposes:} phylogenetics is the weakest case among the
six for quantum-facing benchmarking. The likelihood function is expensive to evaluate
(requires integration over branch lengths), there is no polynomial verifier for
near-optimal topologies, and QUBO encodings of tree topology search remain
speculative for all but the smallest instances.

\subsection*{Heuristic Landscape and Improvement Opportunities}

State-of-the-art tools such as RAxML-NG and IQ-TREE2 implement hill-climbing heuristics 
with SPR (subtree pruning and regrafting) and NNI (nearest-neighbour interchange) moves, 
augmented by random restarts to escape local optima. They are well-engineered and 
practically very competitive for moderate datasets, but face two distinct bottlenecks: 
(i) the likelihood landscape is highly multimodal with many shallow local optima 
separated by long valleys, and (ii) for very large datasets (tens of thousands of taxa), 
per-move likelihood recomputation becomes prohibitively expensive.

AI-guided heuristics are an active and growing area. Azouri et al.\ have shown that 
machine learning can rank candidate SPR moves before their likelihood is computed, 
reducing the number of full likelihood evaluations and speeding up search without 
significant loss of accuracy \cite{Azouri2021}. The same group subsequently demonstrated 
a full reinforcement learning (RL) agent (deep Q-learning) that learns a search 
policy over tree topology space, avoiding greedy uphill moves and reaching 
competitive log-likelihood scores on datasets of up to 20 taxa \cite{AzouriRL2024}.
More recently, end-to-end deep learning approaches such as NeuralNJ combine 
transformer-based distance estimation with RL-guided tree assembly to handle 
hundreds of taxa with improved efficiency \cite{NeuralNJ2025}.
At the same time, purely supervised approaches using CNNs on alignment images or 
site-pattern frequency vectors can predict quartet topologies at the accuracy of 
maximum likelihood, at a fraction of the computational cost---useful as a warm-start 
strategy for larger tree searches \cite{Suvorov2020}.
Quantum algorithms for tree topology search face a fundamental obstacle in encoding 
the discrete and non-metric structure of tree space; no competitive quantum approach 
exists yet.

\subsection*{References}

\begin{itemize}[noitemsep]
  \item \cite{Felsenstein1981}: Foundational ML framework for phylogenetic trees 
    from DNA sequences (J.\ Mol.\ Evol., 1981). \doi{10.1007/BF01734359}
  \item \cite{Yang2014}: Textbook covering substitution models and Bayesian inference 
    (Oxford University Press, 2014). \doi{10.1093/acprof:oso/9780199602605.001.0001}
  \item \cite{Azouri2021}: ML-guided ranking of SPR moves to speed up tree search 
    (Nat.\ Commun., 2021). \doi{10.1038/s41467-021-24761-7}
  \item \cite{AzouriRL2024}: Reinforcement learning agent for phylogenetic tree 
    reconstruction (Mol.\ Biol.\ Evol., 2024). \doi{10.1093/molbev/msae105}
  \item \cite{NeuralNJ2025}: End-to-end deep learning for phylogenetic inference with 
    RL-based tree search (Mol.\ Biol.\ Evol., 2025). \doi{10.1093/molbev/msaf260}
  \item \cite{Suvorov2020}: CNN-based quartet topology inference at ML accuracy 
    (Syst.\ Biol., 2020). \doi{10.1093/sysbio/syz060}
  \item \cite{TreeBASE}: TreeBASE, a database of phylogenetic knowledge
    (Systematic Biology, 1994). \doi{10.1093/sysbio/43.3.514}
  \item \cite{Robinson1981}: Comparison of weighted labelled trees
    (Comb.\ Math.\ \& its Appl., 1981). \doi{10.1007/BFb0089504}
  \item Wikipedia: \url{https://en.wikipedia.org/wiki/Computational_phylogenetics}
\end{itemize}

%%%%%%%%%%%%%%%%%%%%%%%%%%%%%%%%%%%%%%%%%%
\section{Haplotype Phasing (Minimum Error Correction)}
%%%%%%%%%%%%%%%%%%%%%%%%%%%%%%%%%%%%%%%%%%

\subsection*{Problem Description}

Given sequencing reads overlapping heterozygous SNPs, reconstruct the two haplotypes
minimizing inconsistencies.

\subsection*{Mathematical Formulation}

Let $M \in \{-1,0,1\}^{m \times n}$ represent reads $\times$ SNPs
(with $0$ indicating no coverage).

\[
\min_{H \in \{-1,1\}^{2 \times n}}
\sum_{i=1}^{m} \min_{h \in H}
\text{distance}(M_i, h)
\]

where distance counts mismatches over covered SNPs.

\subsection*{Existing Benchmark Datasets}

Haplotype phasing has a relatively well-developed benchmarking ecosystem, centred on
samples with experimentally validated ground-truth haplotypes.

\textbf{Genome in a Bottle (GIAB)} \cite{GIAB2014} provides high-confidence variant
calls and phased haplotypes for a small set of well-characterised human genomes (HG001
through HG007, the Ashkenazi and Chinese trios plus the NIST reference sample NA12878).
The GIAB truth sets are produced by integrating multiple orthogonal sequencing
technologies and are the standard ground truth for benchmarking phasing accuracy
in clinical genomics. Available at \url{https://www.nist.gov/programs-projects/genome-bottle}.

\textbf{Human Pangenome Reference Consortium (HPRC)} \cite{HPRC2023} has released fully
phased, chromosome-scale assemblies of 47 individuals, providing a large panel of
haplotype-resolved genomes at unprecedented completeness. The HPRC dataset represents
the current state-of-the-art for long-read benchmarking at whole-genome scale and is
particularly relevant for testing methods on structurally complex regions.

\textbf{1000 Genomes Project} \cite{1000Genomes2015} provides population-scale genotype
data (2{,}504 individuals, 26 populations) with partially phased haplotypes derived from
statistical phasing. It is the canonical resource for benchmarking population-level
phasing methods such as SHAPEIT and EAGLE.

\textbf{PrecisionFDA Truth Challenge} datasets \cite{PrecisionFDA} have been used in
community benchmarking exercises that include phasing accuracy as an evaluation component,
providing standardised instances and comparison infrastructure.

\textbf{Key advantage for BENQODHI purposes:} haplotype phasing via Minimum Error
Correction (MEC) has a clean NP-hard combinatorial formulation that maps naturally to
QUBO/Ising models. The GIAB and HPRC datasets provide genuine instances with
controllable hardness parameters (read depth, read length, heterozygosity), making this
one of the most tractable problems for quantum-facing benchmark construction.

\subsection*{Heuristic Landscape and Improvement Opportunities}

For short-read data, the field is relatively mature: HapCUT2 (max-cut graph heuristic) 
and WhatsHap (fixed-parameter tractable dynamic programming, optimal in read-length) 
together cover the dominant use cases and produce high-quality haplotype blocks. 
For these settings, improvements in the quality of the MEC solution itself are 
modest; the bottleneck has shifted to \emph{phase block contiguity} and 
\emph{scale} for population-level phasing.

The genuine open frontier is \emph{long-read phasing}. PacBio HiFi and Oxford 
Nanopore reads are long enough to span many heterozygous SNPs, but their higher 
per-base error rate and structural complexity (indels, SVs) break the assumptions 
of short-read MEC solvers. LongPhase addresses this by co-phasing SNPs and SVs 
simultaneously, achieving much larger haplotype blocks ($N_{50} \sim 25$~Mbp vs 
$10$--$15$~Mbp for HapCUT2/WhatsHap) while running an order of magnitude faster 
\cite{LongPhase2022}. MethPhaser further extends phasing contiguity by leveraging 
haplotype-specific methylation signals to bridge SNV-based phase block gaps, 
demonstrating improvements in medically relevant loci including the HLA region 
\cite{MethPhaser2024}.

In terms of AI contributions, deep learning and probabilistic graphical model 
approaches have been proposed for integrating multi-modal signals (reads, methylation, 
population reference panels) into phasing, but remain less mature than the 
graph-cut and DP heuristics. The MEC landscape is highly sparse and structured---a 
property that could in principle be exploited by tensor decomposition or 
message-passing methods, but no dominant approach has emerged yet. Quantum 
annealing for MEC has been formulated as a QUBO but is not yet competitive at 
human-genome scale.

\subsection*{References}

\begin{itemize}[noitemsep]
  \item \cite{Lancia2001}: Single Individual SNP Haplotyping problem and MFR/MSR 
    formulations (ESA, LNCS 2161, 2001). \doi{10.1007/3-540-44676-1\_15}
  \item \cite{BansalBafna2008}: HapCUT max-cut heuristic for haplotype assembly 
    (Bioinformatics, 2008). \doi{10.1093/bioinformatics/btn298}
  \item \cite{LongPhase2022}: LongPhase, ultra-fast co-phasing of SNPs and SVs with 
    long reads (Bioinformatics, 2022). \doi{10.1093/bioinformatics/btac058}
  \item \cite{MethPhaser2024}: MethPhaser, methylation-based phase block chaining 
    (Nat.\ Commun., 2024). \doi{10.1038/s41467-024-49588-0}
  \item \cite{GIAB2014}: Genome in a Bottle Consortium, high-confidence variant
    calls for benchmarking (Nat.\ Biotechnol., 2014). \doi{10.1038/nbt.2835}
  \item \cite{HPRC2023}: Human Pangenome Reference Consortium, draft human pangenome
    reference (Nature, 2023). \doi{10.1038/s41586-023-05896-x}
  \item Wikipedia: \url{https://en.wikipedia.org/wiki/Haplotype}
\end{itemize}

%%%%%%%%%%%%%%%%%%%%%%%%%%%%%%%%%%%%%%%%%%
\section{RNA Secondary Structure Prediction}
%%%%%%%%%%%%%%%%%%%%%%%%%%%%%%%%%%%%%%%%%%

\subsection*{Problem Description}

Predict RNA base pairing structure minimizing free energy subject to structural
constraints.

\subsection*{Mathematical Formulation}

Let $\mathcal{S}$ denote the set of valid pseudoknot-free secondary structures.

\[
\min_{S \in \mathcal{S}} E(S)
\]

Subject to:
\begin{itemize}[noitemsep]
    \item Each base pairs at most once
    \item No crossing pairs (for the standard model)
\end{itemize}

\[
E(S) = \sum_{(i,j)\in S} e_{ij} + \text{loop energies}
\]

where loop energy terms account for hairpin loops, internal loops, and multiloops
according to empirically derived thermodynamic parameters.

\subsection*{Existing Benchmark Datasets}

RNA secondary structure has a well-developed set of benchmark databases, all anchored
on experimentally determined structures from crystallography or chemical probing.

\textbf{bpRNA-1m} \cite{Danaee2018} is the largest and most comprehensive RNA secondary
structure meta-database, containing over 102{,}000 single-molecule known secondary
structures extracted and annotated from seven source databases. Developed by D.~Hendrix's
group at Oregon State University, it provides detailed structural annotations including
stems, loops, bulges, and pseudoknots via the \texttt{bpRNA} annotation tool. The
standard machine learning split (bpRNA-1m(90), removing sequences $\ge$90\% identical)
contains over 28{,}000 non-redundant entries and has become the default training/testing
set for deep learning methods. Available at \url{http://bprna.cgrb.oregonstate.edu}.

\textbf{ArchiveII} \cite{Mathews2004} is the standard small-scale benchmark, containing
3{,}857 structures from nine well-characterised RNA families (tRNA, 5S rRNA, SRP,
RNaseP, tmRNA, Group I introns, telomerase RNA, 16S rRNA, 23S rRNA). It has been
used continuously for over 20 years as a cross-method comparison set.

\textbf{Rfam} \cite{Rfam} is the comprehensive RNA family database, currently containing
$>$4{,}000 families with consensus secondary structures validated by covariance models.
The \textbf{bpRNA-new} dataset \cite{MXfold22021}, derived from Rfam 14.4 but excluding
families present in Rfam 12.2, is the canonical out-of-distribution test set for
assessing generalisation of deep learning models to unseen RNA families.

\textbf{PDB-RNA} provides experimentally determined structures from the Protein Data Bank,
serving as an independent evaluation set for methods trained on Rfam-derived data.

\textbf{Eterna100} \cite{Wayment2022} is a curated set of 100 challenging RNA design
puzzles from the Eterna crowdsourcing platform, covering a wide range of structural
motifs and providing a stringent test of thermodynamic model accuracy.

\textbf{Key limitation for BENQODHI purposes:} for non-pseudoknotted structures, the
MFE problem is solvable in $O(n^3)$ by classical dynamic programming and is therefore
not a hard optimisation target for quantum approaches. The relevant hard variant is
pseudoknot-containing structures, for which the problem becomes NP-hard; however the
thermodynamic model for pseudoknots is less well-validated, which complicates
benchmark design.

\subsection*{Heuristic Landscape and Improvement Opportunities}

This is arguably the problem with the highest potential for AI-driven improvement among 
the six, as the classical DP approach (ViennaRNA/mfold) has remained essentially 
unchanged for over four decades and has hit an accuracy ceiling of around 80\% F1 
on standard benchmarks, largely due to the imprecision of the empirical energy 
parameter sets.

Deep learning has genuinely disrupted this landscape. SPOT-RNA was among the first 
to demonstrate that ensembles of two-dimensional convolutional networks trained with 
transfer learning could predict pseudoknot base pairs---a class of structures that 
DP cannot handle---significantly exceeding classical methods on noncanonical pairs 
\cite{SPOTRNA2019}. UFold extended this with U-Net architectures operating on 
base-pair probability matrices, reaching F1 $= 0.91$ on within-family benchmarks 
compared to $0.55$--$0.84$ for classical methods \cite{UFold2022}. 
EternaFold, trained on over 20,000 synthetic RNA constructs from the Eterna 
crowdsourcing platform, consistently outperforms thermodynamic models on viral 
genomes, mRNAs, and riboswitches, and is directly relevant to mRNA vaccine design 
\cite{Wayment2022}.
KnotFold addresses the pseudoknot problem using a learned potential combined 
with a minimum-cost flow algorithm that lifts the nested structure constraint, 
enabling full pseudoknot prediction at competitive accuracy \cite{KnotFold2024}.

A critical and underappreciated caveat, however, is \emph{generalisation across 
RNA families}: Szikszai et al.\ showed that most deep learning models achieve 
impressive intra-family metrics but fail on out-of-distribution families, 
suggesting they overfit structural patterns rather than learning fundamental 
folding physics \cite{Szikszai2022}. This points to hybrid thermodynamic--ML models 
as the most promising direction, and cross-family benchmarking as an essential 
evaluation criterion. Quantum approaches have been proposed for small RNA folding 
instances via QUBO mappings but are not yet at any practical scale.

\subsection*{References}

\begin{itemize}[noitemsep]
  \item \cite{ZukerStiegler1981}: Original DP algorithm for MFE RNA folding 
    (Nucleic Acids Res., 1981). \doi{10.1093/nar/9.1.133}
  \item \cite{Mathews1999}: Expanded thermodynamic parameters for RNA structure 
    prediction (J.\ Mol.\ Biol., 1999). \doi{10.1006/jmbi.1999.2700}
  \item \cite{Danaee2018}: bpRNA, large-scale automated annotation and analysis of RNA
    secondary structure (Nucleic Acids Res., 2018). \doi{10.1093/nar/gky285}
  \item \cite{Rfam}: Rfam, the RNA family database (Nucleic Acids Res., 2021).
    \doi{10.1093/nar/gkaa1047}
  \item \cite{SPOTRNA2019}: SPOT-RNA, deep learning with transfer learning for 
    pseudoknot prediction (Nat.\ Commun., 2019). \doi{10.1038/s41467-019-13395-9}
  \item \cite{UFold2022}: UFold, U-Net architecture for fast RNA structure prediction 
    (Nucleic Acids Res., 2022). \doi{10.1093/nar/gkab1074}
  \item \cite{Wayment2022}: EternaFold, multitask learning on crowdsourced RNA data 
    (Nature Methods, 2022). \doi{10.1038/s41592-022-01605-0}
  \item \cite{KnotFold2024}: KnotFold, learned potentials with minimum-cost flow for 
    pseudoknot prediction (Commun.\ Biol., 2024). \doi{10.1038/s42003-024-05952-4}
  \item \cite{Szikszai2022}: Critical evaluation of DL generalisation across RNA families 
    (Bioinformatics, 2022). \doi{10.1093/bioinformatics/btac415}
  \item Wikipedia: \url{https://en.wikipedia.org/wiki/Nucleic_acid_secondary_structure}
\end{itemize}

%%%%%%%%%%%%%%%%%%%%%%%%%%%%%%%%%%%%%%%%%%
\section{Radiotherapy Treatment Planning}
%%%%%%%%%%%%%%%%%%%%%%%%%%%%%%%%%%%%%%%%%%

\subsection*{Problem Description}

Determine radiation beam intensities to deliver prescribed dose to a tumour while
sparing healthy tissue and organs at risk.

The standard formulation is a constrained quadratic programme; extensions
incorporating biological objectives (TCP, NTCP) become nonconvex.

\subsection*{Mathematical Formulation}

Let $x \ge 0$ be beam intensities (fluence map), $A$ the dose-influence matrix,
and $d$ the prescribed dose vector.

\[
\min_{x \ge 0}
\| A x - d \|^2 + \lambda R(x)
\]

Subject to:

\[
(Ax)_j \le d^{\text{max}}_j \quad \text{for organs at risk } j
\]

where $R(x)$ is a regularisation term (e.g.\ smoothness penalty on the fluence map).

\subsection*{Existing Benchmark Datasets}

Radiotherapy treatment planning is the best-served of the six problems in terms of
open benchmark infrastructure explicitly designed for algorithmic comparison.

\textbf{CORT (Common Optimisation for Radiation Therapy)} \cite{Craft2014} is the primary
open benchmark for IMRT optimisation research, created by D.~Craft's group at Harvard/MGH.
It provides pre-computed dose-influence matrices for four clinical cases (a TG-119 phantom,
prostate, liver, and head-and-neck) at multiple beam angle configurations, together with
suggested clinical constraints. The dose-influence matrix is the core object needed to
formulate and solve the fluence optimisation problem, so CORT allows researchers to
compare optimisation algorithms directly without requiring radiation physics infrastructure.
Available at \url{https://github.com/rcruces/cort_dataset}.

\textbf{AAPM TG-119} \cite{Ezzell2009} is the standard clinical commissioning phantom
defined by the American Association of Physicists in Medicine. It defines a C-shaped
target wrapping around an organ at risk, creating a genuinely difficult dose-shaping
problem, and has served as a de facto benchmark across dozens of algorithmic studies.
The TG-119 case is included in the CORT dataset.

\textbf{TROTS (The Radiotherapy Optimisation Test Set)} \cite{Breedveld2017} is a larger
collection providing clinical cases for multiple tumour sites (head and neck, prostate,
rectum, breast, lung, brain) with clinical planning objectives and constraints,
intended to enable systematic comparison of multi-criteria optimisation approaches.

\textbf{OpenKBP} \cite{OpenKBP2021} provides 340 head-and-neck patient cases with
full CT images, structure contours, and dose distributions, designed specifically for
knowledge-based planning and dose-prediction benchmarks. While more focused on
dose prediction than fluence optimisation, it complements CORT for learning-based
approaches.

\textbf{Key advantage for BENQODHI purposes:} the CORT dataset provides directly
usable optimisation instances with known problem structure (convex QP for fluence
optimisation, combinatorial for beam angle selection). Classical solver baselines are
well-characterised, making approximation ratios computable. This is arguably the
most mature benchmark infrastructure among the six problems for quantum-facing
repurposing, requiring primarily the addition of QUBO reformulations of the beam
angle selection sub-problem.

\subsection*{Heuristic Landscape and Improvement Opportunities}

The standard convex IMRT problem is well-solved in clinical practice by commercial 
systems (Eclipse, RayStation) using gradient-based interior-point methods. 
The real bottlenecks lie in three extensions where current heuristics leave 
clear room for improvement.

First, \emph{beam angle optimisation} (BAO) is a combinatorial problem that standard 
systems address with coarse fixed-angle heuristics or exhaustive evaluation of a 
small candidate set. Deep learning classifiers trained to mimic column-generation 
solvers can predict near-optimal beam configurations in seconds rather than hours 
\cite{Barkousaraie2020}. Reinforcement learning agents have been applied to the same 
problem, learning to sequentially select beam angles using a Monte Carlo tree search 
with a graph neural network policy, achieving plans comparable to or better than the 
RL-free DL baseline \cite{Bao2023}.

Second, \emph{online adaptive replanning}---rescaling dose on each fraction to account 
for daily anatomical changes---requires new plans in minutes rather than hours, 
making classical iterative optimisation impractical. Conditional generative adversarial 
networks (cGANs) can predict deliverable fluence maps directly from updated anatomy 
\cite{Buchanan2023}, while deep RL agents trained in a simulated environment 
(multi-objective adjustment policy networks, MOAPN) have been shown to adjust 
treatment planning objectives autonomously within commercial planning systems \cite{MuAPNLung2023}.

Third, incorporating \emph{biological objectives} (tumour control probability, normal 
tissue complication probability) introduces non-linear structure that standard 
dose-volume heuristics handle imprecisely. Quantum deep RL has been proposed for 
adaptive radiotherapy under uncertainty, using quantum states to model biological 
stochasticity and reporting improved clinical decision quality on NSCLC cohorts 
\cite{NiraulaQDRL2021}.

\subsection*{References}

\begin{itemize}[noitemsep]
  \item \cite{Bortfeld2006}: Review of IMRT optimisation (Phys.\ Med.\ Biol., 2006).
    \doi{10.1088/0031-9155/51/13/R21}
  \item \cite{Craft2007}: Multi-criteria IMRT optimisation (IJROBP, 2007).
    \doi{10.1016/j.ijrobp.2007.08.019}
  \item \cite{Craft2014}: CORT dataset, shared data for IMRT optimisation research
    (GigaScience, 2014). \doi{10.1186/2047-217X-3-37}
  \item \cite{Ezzell2009}: AAPM TG-119 report, IMRT commissioning benchmark
    (Med.\ Phys., 2009). \doi{10.1118/1.3238104}
  \item \cite{Breedveld2017}: TROTS, the radiotherapy optimisation test set
    (Data in Brief, 2017). \doi{10.1016/j.dib.2017.03.019}
  \item \cite{OpenKBP2021}: OpenKBP, open-access knowledge-based planning dataset
    (Med.\ Phys., 2021). \doi{10.1002/mp.14845}
  \item \cite{Barkousaraie2020}: DL-based beam angle optimisation mimicking column 
    generation (Med.\ Phys., 2020). \doi{10.1002/mp.13986}
  \item \cite{Bao2023}: RL with Monte Carlo tree search for beam angle optimisation 
    (arXiv, 2023). \url{https://arxiv.org/abs/2302.12703}
  \item \cite{Buchanan2023}: cGAN for adaptive plan prediction in MR-guided RT 
    (Front.\ Oncol., 2023). \doi{10.3389/fonc.2022.939951}
  \item \cite{MuAPNLung2023}: Multi-agent deep RL for automated IMRT planning 
    (Front.\ Oncol., 2023). \doi{10.3389/fonc.2023.1124458}
  \item \cite{NiraulaQDRL2021}: Quantum deep RL for adaptive radiotherapy 
    (Sci.\ Rep., 2021). \doi{10.1038/s41598-021-02910-y}
  \item Wikipedia: \url{https://en.wikipedia.org/wiki/Intensity-modulated_radiation_therapy}
\end{itemize}

%%%%%%%%%%%%%%%%%%%%%%%%%%%%%%%%%%%%%%%%%%
\section{Metabolic Flux Optimisation (Flux Balance Analysis)}
%%%%%%%%%%%%%%%%%%%%%%%%%%%%%%%%%%%%%%%%%%

\subsection*{Problem Description}

Predict metabolic fluxes satisfying stoichiometric constraints while optimizing a 
biological objective (e.g.\ growth rate, ATP production).

The base FBA is a linear programme; extensions (mixed-integer FBA, dynamic FBA,
thermodynamic FBA) introduce combinatorial or nonconvex structure.

\subsection*{Mathematical Formulation}

Let $S \in \mathbb{R}^{m \times n}$ be the stoichiometric matrix and 
$v \in \mathbb{R}^n$ the flux vector.

\[
\max_{v} \; c^T v
\]

Subject to:

\[
S v = 0, \qquad v_{\min} \le v \le v_{\max}
\]

This is a linear programme; genome-scale models ($n \sim 10^3$--$10^4$) and 
multi-objective or mixed-integer variants present significant challenges.

\subsection*{Existing Benchmark Datasets}

FBA has the most developed and standardised benchmark infrastructure of the six problems,
including an explicit solver comparison study of direct relevance to BENQODHI.

\textbf{BiGG Models} \cite{BiGG2016,BiGG2020} is the canonical repository for
high-quality, manually curated genome-scale metabolic models, currently containing over
100 reconstructions for organisms spanning bacteria, yeast, and human. Models are
standardised with consistent metabolite and reaction identifiers, and quality is assessed
using the \textbf{MEMOTE} scoring system \cite{MEMOTE2020}. Key models include:
\textit{iJO1366} (E.~coli, 1{,}366 genes, 2{,}583 reactions), \textit{iYali4}
(Yarrowia lipolytica), \textit{Recon3D} (human, 13{,}543 reactions, 4{,}140 metabolites).
Available at \url{http://bigg.ucsd.edu}.

\textbf{EMBL GEM collection} \cite{Heinken2023} provides automatically reconstructed
genome-scale models for 7{,}302 human microbiome organisms, enabling community-level FBA
simulations. Used for large-scale solver benchmarking at community scale.

\textbf{Solver benchmarks on BiGG/EMBL data:} the most directly relevant existing
benchmark for BENQODHI is the 2024 study by D.~Machado (NTNU) \cite{Machado2024}, which
systematically compares LP and MILP solvers (Gurobi, CPLEX, SCIP, HiGHS, GLPK, COIN)
across models of increasing size, covering both single-species FBA and community
simulations. This study provides reproducible baseline runtimes across problem sizes up
to Recon3D ($\sim$10K variables) and communities of up to 20 species, and is the
natural starting point for extending towards quantum solvers.

\textbf{$^{13}$C flux measurement datasets} provide ground-truth flux distributions
measured experimentally via isotope labelling, used to assess prediction accuracy rather
than solver speed. Key datasets include E.~coli cultures under defined carbon sources
\cite{Antoniewicz2013} and yeast fermentation datasets.

\textbf{Key advantage for BENQODHI purposes:} the LP formulation is the cleanest and
most tractable of the six problems for quantum-facing benchmarking. Solver comparisons
already exist (Machado 2024), the problem is natively structured as an LP/MILP with
well-defined instances of varying size, and QUBO reformulations of the MILP knockout
screening variant are natural. The gap between GLPK/HiGHS (competitive open-source)
and Gurobi (commercial) on Recon3D-scale instances represents a concrete target for
quantum hybrid approaches.

\subsection*{Heuristic Landscape and Improvement Opportunities}

The base FBA LP is trivially solved for any genome-scale model and is not the 
computational bottleneck. The open challenges lie in \emph{model accuracy} 
and \emph{hard extensions}.

On model accuracy: classical FBA requires specification of uptake flux bounds 
from extracellular concentrations, a conversion that involves complex transporter 
kinetics. Hybrid neural-mechanistic models (augmented metabolic networks, AMNs) 
address this directly by learning a neural pre-processing layer that maps 
experimental medium conditions to input flux bounds, systematically outperforming 
classical FBA at predicting growth rates and secretion phenotypes while requiring 
far smaller training sets than pure data-driven approaches \cite{Millard2023}.
Graph neural networks integrated with FBA (FlowGAT) have been shown to improve 
prediction of gene essentiality by exploiting the network structure of metabolic 
fluxes as a graph representation \cite{FlowGAT2024}.

On hard extensions: mixed-integer FBA for optimising combinatorial gene 
knockout strategies requires searching a large discrete space. Standard approaches use 
branch-and-bound MILP solvers, but ML surrogate models trained on FBA 
outputs can accelerate screening of knockout combinations by orders of magnitude 
\cite{Zampieri2019,Sahu2021}. Dynamic FBA involves coupled ODE/LP systems 
with no general-purpose efficient solver; recurrent neural networks have been 
used to emulate dynamic FBA at reduced computational cost. Multi-omics integration 
(transcriptomics, proteomics, metabolomics) into context-specific GEMs remains 
an active methodological frontier where machine learning-based data fusion is 
outperforming classical constraint-based integration methods \cite{GonzalezFlux2023}.

\subsection*{References}

\begin{itemize}[noitemsep]
  \item \cite{VarmaPalsson1994}: Stoichiometric FBA framework for E.~coli 
    (Appl.\ Environ.\ Microbiol., 1994). \doi{10.1128/aem.60.10.3724-3731.1994}
  \item \cite{Orth2010}: Review of FBA methodology and extensions 
    (Nat.\ Biotechnol., 2010). \doi{10.1038/nbt.1614}
  \item \cite{BiGG2016}: BiGG Models, platform for integrating and sharing genome-scale
    models (Nucleic Acids Res., 2016). \doi{10.1093/nar/gkv1049}
  \item \cite{BiGG2020}: BiGG Models 2020, multi-strain models and expanded phylogenetic
    coverage (Nucleic Acids Res., 2020). \doi{10.1093/nar/gkz1055}
  \item \cite{MEMOTE2020}: MEMOTE, standardised genome-scale metabolic model testing
    (Nat.\ Biotechnol., 2020). \doi{10.1038/s41587-020-0446-y}
  \item \cite{Heinken2023}: Genome-scale metabolic reconstruction of 7{,}302 human
    microorganisms (Nat.\ Biotechnol., 2023). \doi{10.1038/s41587-022-01628-0}
  \item \cite{Machado2024}: Benchmark of optimisation solvers for genome-scale metabolic
    modeling (mSystems, 2024). \doi{10.1128/msystems.00833-23}
  \item \cite{Zampieri2019}: ML and deep learning applied to genome-scale metabolic 
    modelling (PLoS Comput.\ Biol., 2019). \doi{10.1371/journal.pcbi.1007084}
  \item \cite{Sahu2021}: Review of advances in FBA via ML and mechanism-based models 
    (Comput.\ Struct.\ Biotechnol.\ J., 2021). \doi{10.1016/j.csbj.2021.08.004}
  \item \cite{Millard2023}: Neural-mechanistic hybrid models for FBA phenotype 
    prediction (Nat.\ Commun., 2023). \doi{10.1038/s41467-023-40380-0}
  \item \cite{FlowGAT2024}: FlowGAT, graph neural network for gene essentiality 
    via FBA integration (npj Syst.\ Biol.\ Appl., 2024). \doi{10.1038/s41540-024-00348-2}
  \item \cite{GonzalezFlux2023}: ML-based flux prediction from omics data 
    (Comput.\ Struct.\ Biotechnol.\ J., 2023). \doi{10.1016/j.csbj.2023.10.002}
  \item Wikipedia: \url{https://en.wikipedia.org/wiki/Flux_balance_analysis}
\end{itemize}

\newpage

\section*{Summary Table}

\begin{center}
\begin{tabular}{llccc}
\toprule
Problem & Key Benchmark Datasets & Type & AI Frontier & Quantum-Facing \\
\midrule
MSA & BAliBASE, HomFam, SABmark & Combinatorial & Transformer aligners & Limited \\
Phylogenetic Tree & TreeBASE, simulated (seq-gen) & Combinatorial & RL tree search & Speculative \\
Haplotype Phasing & GIAB, HPRC, 1000 Genomes & Combinatorial & Long-read multi-modal & Strong \\
RNA Structure & bpRNA-1m, ArchiveII, Rfam & Energy min. & DL energy models & Hard variant only \\
Radiotherapy & CORT, TROTS, TG-119 & QP / Combinatorial & RL + cGAN planning & Strong \\
Metabolic Flux & BiGG, EMBL GEMs, Machado 2024 & LP / MILP & Hybrid neural-mech. & Strong \\
\bottomrule
\end{tabular}
\end{center}

\section*{Conclusion}

These six optimisation problems represent foundational and still computationally
challenging tasks in biology and bioinformatics. For each, the open question is
how much further current heuristics can be pushed---in solution quality, in speed,
and in scalability to modern data regimes.
Classical methods remain strong baselines, but AI-driven approaches (transformers, 
reinforcement learning, graph neural networks, hybrid mechanistic-ML models) are 
demonstrably improving heuristic performance across all six domains. 

Existing benchmark datasets provide a rich starting point: BiGG/EMBL and CORT offer
the most direct raw material for quantum-facing repurposing, as they combine clean
mathematical formulations with reproducible instances and established classical baselines.
BAliBASE, bpRNA-1m, GIAB, and TreeBASE require additional reformulation work to
expose computational hardness in a solver-agnostic way. In all cases, the key step
BENQODHI must take is to augment domain-specific quality metrics with hardness-characterising
metrics---approximation ratios, scaling exponents, and time-to-target-quality---enabling
genuine cross-paradigm comparison between classical, AI-based, and quantum solvers.

\bibliography{references}

\end{document}